\textbf{E01.} El paciente podrá reservar una cita eligiendo especialidad, médico, fecha y hora entre los horarios disponibles.\par
\textbf{E02.} La plataforma enviará al paciente un recordatorio por mensaje de texto 24 horas antes de su cita.\par
\textbf{E03.} La aplicación debe ser rápida.\par
\textbf{E04.} El recepcionista podrá registrar la llegada del paciente que se acerca a la ventanilla.\par
\textbf{E05.} La plataforma se desarrollará en Java 21 con base de datos PostgreSQL, tecnologías que domina el área de TI.\par
\textbf{E06.} El médico podrá consultar la lista de pacientes del día y marcar cada consulta como atendida.\par
\textbf{E07.} Los datos de salud de los pacientes deben estar seguros.\par
\textbf{E08.} La aplicación móvil debe funcionar en teléfonos con Android 10 o superior.\par
\textbf{E09.} Una cita cancelada con menos de cuatro horas de anticipación se registrará como inasistencia.\par
\textbf{E10.} La plataforma debe ser fácil de usar para los pacientes adultos mayores.\par
\textbf{E11.} El proyecto debe entregarse en un plazo máximo de seis meses.\par
\textbf{E12.} La sesión de un usuario se cerrará automáticamente tras 10 minutos sin actividad.\par
\textbf{E13.} El quiosco de la sala de espera leerá el número de cédula ingresado en su pantalla táctil e imprimirá un ticket con el número de turno.\par
\textbf{E14.} La plataforma debe estar siempre disponible.\par
\textbf{E15.} El director médico podrá generar un reporte mensual de inasistencias por especialidad.\par
\textbf{E16.} Cada incremento se validará con el director médico al final de cada iteración de dos semanas.\par
\textbf{E17.} El tratamiento de los datos personales de los pacientes debe cumplir la Ley Orgánica de Protección de Datos Personales del Ecuador.\par
\textbf{E18.} La plataforma debe integrarse sin problemas con el sistema de facturación actual.\par
\textbf{E19.} Si el servidor principal falla, la plataforma debe restablecer el servicio en menos de 30 minutos sin perder ninguna cita confirmada.\par
\textbf{E20.} Todo cambio en el código fuente será revisado por otro integrante del equipo antes de integrarse a la rama principal.\par
